% !TeX root = ../教材题库_全部条目_转写工作稿.tex
% Source: XJ-B1-C05-HW-P162-Y1346
% 第5章 / 5.1 任意角和弧度制 / 教材 p.162 / PDF 第169页
% OCR未校对：公式与文字均可能有误；未完成数学TeX转写；禁止直接发给学生
% 原题及相邻原解截图（核对用）：../题目原图/XJ-B1-C05-HW-P162-Y1346-01.webp
% 本文件只在转写完成后作为题干使用；OCR 条目尚不是可靠题干。
\begin{bankitem}{习题5.1 第13题}{第5章 / 5.1 任意角和弧度制 / 教材 p.162 / PDF 第169页}{XJ-B1-C05-HW-P162-Y1346}{OCR未校对：公式与文字均可能有误；未完成数学TeX转写；禁止直接发给学生}
\begin{ocrdraft}
[教材p.162 / PDF第169页]
\par
13．已知相互咬合的两个齿轮，大轮有48齿，小轮有20齿，当大轮顺时针转
\par
动一周时，小轮转动的角是多少度？多少弧度？如果大轮的转速是150r/mm，小轮
\par
的半径为10cm，那么小轮圆周上的点每秒转过的弧长是多少？
\end{ocrdraft}
\end{bankitem}
